% =====================================================================
\section{FreeRTOS: задачи и синхронизация}\label{les:rtos}
% =====================================================================
\lessonintro
  {разделить программу на независимые задачи, организовать безопасный обмен данными, задействовать оба ядра}
  {уроки 1--7 (неблокирующий \code{loop()}, кнопка, АЦП, экран), прерывания из урока 3}
  {кнопка, экран и Serial работают как отдельные задачи}
  {ничего нового из «железа»}

\subsection{Зачем нужна операционная система реального времени}

В конце урока 7 \code{loop()} совмещал пять дел, и каждое новое дело замедляло остальные.
FreeRTOS позволяет написать каждое дело как отдельную «программу» --- \textbf{задачу} --- со своим
бесконечным циклом и приоритетом. Планировщик переключает задачи каждую миллисекунду, а на
ESP32-S3 задачи могут выполняться на двух ядрах по-настоящему параллельно.

Вы уже пользовались FreeRTOS, не зная об этом: \code{setup()} и \code{loop()} выполняются внутри
задачи \code{loopTask} на ядре~1, а \code{delay()} --- это вызов FreeRTOS \code{vTaskDelay()}, который
отдаёт процессор другим задачам.

\begin{figure}[H]
\centering
\begin{tikzpicture}[state/.style={blk,circle,minimum size=19mm,fill=blkmem}]
  \node[state,fill=blkrf] (run) at (0,0) {Выпол-\\няется};
  \node[state] (ready) at (-5,0) {Готова};
  \node[state,fill=blkper] (blk) at (5,0) {Блоки-\\рована};
  \node[state,fill=blkrtc] (susp) at (0,-4.2) {Приоста-\\новлена};
  \draw[arr] (ready) to[bend left=25] node[above,font=\scriptsize]{выбрана планировщиком} (run);
  \draw[arr] (run) to[bend left=25] node[below,font=\scriptsize]{вытеснена} (ready);
  \draw[arr] (run) to[bend left=25] node[above,font=\scriptsize,align=center]{\code{vTaskDelay}, ждёт\\ очередь/мьютекс} (blk);
  \draw[arr] (blk.south) .. controls (3.5,-2.1) and (-3.5,-2.1) .. node[below,font=\scriptsize,pos=0.3]{событие / таймаут} (ready.south);
  \draw[arr] (run) -- node[right,font=\scriptsize,pos=0.75]{\code{vTaskSuspend}} (susp);
  \draw[arr] (susp.west) to[bend left=25] node[pos=0.35,below left,font=\scriptsize]{\code{vTaskResume}} (ready.south west);
\end{tikzpicture}
\caption{Состояния задачи FreeRTOS. Задача, ждущая в \code{vTaskDelay()}, не тратит процессорное время}
\end{figure}

\subsection{Первые задачи}

\begin{figure}[H]
\centering
\begin{tikzpicture}[x=0.9cm,y=0.9cm]
  \node[left,font=\sffamily\small] at (0,1.5) {Ядро 0};
  \node[left,font=\sffamily\small] at (0,0) {Ядро 1};
  \fill[gray!12] (0,1.2) rectangle (14.5,1.8);
  \fill[gray!12] (0,-0.3) rectangle (14.5,0.3);
  \node[font=\scriptsize,text=gray] at (13.5,1.5) {idle};
  \node[font=\scriptsize,text=gray] at (13.5,0) {idle};
  \foreach \s in {0,3,6,9,12}{
    \draw[fill=blkmem] (\s,1.2) rectangle ++(0.45,0.6);
    \node[above,font=\scriptsize] at (\s+0.22,1.8) {b1};}
  \foreach \s in {0.1,7.1}{
    \draw[fill=blkmem] (\s,-0.3) rectangle ++(0.45,0.6);
    \node[below,font=\scriptsize] at (\s+0.22,-0.3) {b2};}
  \draw[fill=blkcpu] (0.7,-0.3) rectangle ++(0.9,0.6) node[midway,font=\scriptsize]{loop};
  \draw[->] (0,-1.2) -- (15,-1.2) node[right,font=\scriptsize]{$t$};
  \foreach \x/\l in {0/0,3/300,7/700,14/1400}
    \draw (\x,-1.15) -- (\x,-1.25) node[below,font=\scriptsize]{\l\,мс};
\end{tikzpicture}
\caption{Задачи из листинга ниже: \code{b1} (GPIO4, каждые 300 мс) на ядре 0, \code{b2} (GPIO7, каждые 700 мс)
и \code{loop()} на ядре 1. Каждая работает доли миллисекунды, остальное время ядро свободно (idle)}
\end{figure}

\codecap{Две независимые задачи (отдельный эксперимент, светодиоды на GPIO4 и GPIO7)}
\codeinput{l7_two_tasks}

Сравните с уроком 1: там для двух частот мигания пришлось вести отдельные таймеры на \code{millis()},
а здесь каждая задача просто «спит» столько, сколько ей нужно.

\subsection{Общие данные: очередь и мьютекс}

Когда несколько задач работают с одними данными, возникает опасность: одна задача может прочитать
структуру в тот момент, когда другая записала её лишь наполовину. Для передачи данных между
задачами используется \textbf{очередь} --- потокобезопасный FIFO-буфер.

\begin{figure}[H]
\centering
\begin{tikzpicture}
  \node[blk,fill=blkper,minimum width=3cm,minimum height=1.2cm] (p) {Задача-производитель\\ \scriptsize читает АЦП};
  \foreach \i/\v in {0/1820,1/1835,2/1790,3/,4/}{
    \node[draw,minimum width=9mm,minimum height=8mm,fill=\ifx\v\empty white\else blkmem\fi,font=\ttfamily\scriptsize]
      (c\i) at ($(p.east)+(2.4+\i*0.9,0)$) {\v};
  }
  \node[blk,fill=blkcpu,minimum width=2.9cm,minimum height=1.2cm,right=29mm of c4] (c) {Задача-потребитель\\ \scriptsize печатает в Serial};
  \draw[arr] (p) -- node[above,font=\ttfamily\scriptsize]{xQueueSend} (c0);
  \draw[arr] (c4) -- node[above,font=\ttfamily\scriptsize]{xQueueReceive} (c);
  \node[font=\sffamily\scriptsize,below=1mm of c2] {очередь на 5 элементов};
\end{tikzpicture}
\caption{Очередь развязывает задачи по времени}
\end{figure}

\codecap{Производитель и потребитель}
\codeinput{l7_queue}

Если же данные --- это общее \emph{состояние}, которое все читают и меняют (например, «светодиод включён, яркость 60 \%»),
удобнее \textbf{мьютекс}: «ключ», который задача берёт перед работой с данными и возвращает после.
Пока ключ у одной задачи, остальные ждут.

\begin{table}[H]
\centering\small
\begin{tabularx}{\textwidth}{@{}l X@{}}
\toprule
Примитив & Когда использовать \\
\midrule
Очередь (\code{xQueue*}) & передача данных между задачами или из прерывания (\code{xQueueSendFromISR}) \\
Мьютекс (\code{xSemaphoreCreateMutex}) & защита общего ресурса: переменных, \code{Serial}, шины I\textsuperscript{2}C \\
Двоичный семафор & сигнал «событие произошло» (например, из прерывания в задачу) \\
Уведомления (\code{xTaskNotify}) & самый быстрый и лёгкий сигнал конкретной задаче \\
Программный таймер (\code{xTimerCreate}) & периодический вызов короткой функции без отдельной задачи \\
\bottomrule
\end{tabularx}
\caption{Основные средства синхронизации FreeRTOS}
\end{table}

\subsection{Практика: мьютекс --- «ключ» от общего ресурса}

Когда две задачи пользуются одной вещью --- например, \code{Serial}, --- они могут мешать друг другу. Проверим!
Две задачи печатают свои фразы по одной букве. Сначала поставь в программе \code{USE\_MUTEX = false}:

\codecap{Две задачи и мьютекс}
\codeinput{l13_mutex}

Без мьютекса в Serial Monitor получится каша: «Кошка говСобака говорит: гаорит: мяу\ldots». Задачи
перебивают друг друга на полуслове. Теперь поставь \code{true}: фразы идут ровно, по очереди.

\textbf{Мьютекс} (от англ. \emph{mutual exclusion} --- взаимное исключение) работает как ключ от единственной
комнаты: задача берёт ключ (\code{xSemaphoreTake}), делает свою работу и возвращает ключ
(\code{xSemaphoreGive}). Пока ключ у одной задачи, другая ждёт.

\begin{figure}[H]
\centering
\begin{tikzpicture}[x=1cm,y=1cm]
  \node[font=\sffamily\small,anchor=east] at (-0.2,1) {без мьютекса};
  \foreach \s/\e/\c in {0/0.5/esp,0.5/1/accent,1/1.5/esp,1.5/2/accent,2/2.5/esp,2.5/3/accent,3/3.5/esp,3.5/4/accent,4/4.5/esp,4.5/5/accent}
    \fill[\c!40] (\s,0.8) rectangle (\e,1.2);
  \node[font=\sffamily\scriptsize,anchor=west] at (5.2,1) {буквы вперемешку};
  \node[font=\sffamily\small,anchor=east] at (-0.2,0) {с мьютексом};
  \fill[esp!40] (0,-0.2) rectangle (2.5,0.2) node[midway,font=\sffamily\scriptsize]{кошка};
  \fill[accent!40] (2.5,-0.2) rectangle (5,0.2) node[midway,font=\sffamily\scriptsize]{собака};
  \node[font=\sffamily\scriptsize,anchor=west] at (5.2,0) {фразы целиком, по очереди};
\end{tikzpicture}
\caption{Мьютекс не даёт задачам перебивать друг друга}
\end{figure}

\begin{table}[H]
\centering\small
\begin{tabularx}{\textwidth}{@{}l X@{}}
\toprule
Инструмент FreeRTOS & Когда пригодится \\
\midrule
Задача (\code{xTaskCreate}) & отдельное дело со своим циклом \\
Очередь (\code{xQueue\ldots}) & передать данные из одной задачи в другую (или из прерывания) \\
Мьютекс (\code{xSemaphoreCreateMutex}) & не дать двум задачам одновременно пользоваться \code{Serial}, экраном, общей переменной \\
Уведомление (\code{xTaskNotify}) & разбудить конкретную задачу: «событие произошло» \\
Программный таймер (\code{xTimerCreate}) & вызывать короткую функцию периодически \\
\bottomrule
\end{tabularx}
\caption{Главные инструменты FreeRTOS}
\end{table}

\begin{caution}
Если задаче не хватает памяти стека, плата перезагружается с сообщением \code{Stack canary watchpoint triggered}.
Тогда увеличь третий параметр \code{xTaskCreatePinnedToCore()} --- например, с 2048 до 4096.
\end{caution}


\subsection{Самостоятельная работа: устройство на задачах}

\begin{figure}[H]
\centering
\begin{tikzpicture}[node distance=8mm]
  \node[blk,fill=blkmem,minimum width=34mm,minimum height=14mm] (st) {\code{state}\\ \scriptsize + мьютекс \code{stateMutex}};
  \node[blk,fill=blkper,minimum width=30mm,minimum height=12mm,left=22mm of st] (in) {\code{inputTask}\\ \scriptsize кнопка, потенциометр · 10 мс\\ \scriptsize приоритет 2};
  \node[blk,fill=blkrf,minimum width=30mm,minimum height=12mm,right=22mm of st] (dsp) {\code{displayTask}\\ \scriptsize экран · 100 мс\\ \scriptsize приоритет 1};
  \node[blk,fill=blkcpu,minimum width=30mm,minimum height=12mm,below=10mm of st] (lp) {\code{loop()}\\ \scriptsize команды Serial};
  \draw[arr] (in) -- node[above,font=\scriptsize]{\code{toggle()}} node[below,font=\scriptsize]{\code{setLevel()}} (st);
  \draw[arr] (st) -- node[above,font=\scriptsize]{\code{getState()}} (dsp);
  \draw[darr] (lp) -- node[right,font=\scriptsize]{\code{setOn()}, \code{getState()}} (st);
  \node[below=1mm of dsp,font=\sffamily\scriptsize,align=center] {единственная задача,\\ работающая с I\textsuperscript{2}C};
\end{tikzpicture}
\caption{Устройство на задачах: задачи обращаются к общему состоянию только через функции под мьютексом}
\end{figure}

Главное правило: общее состояние меняется только через функции \code{toggle()}, \code{setLevel()},
\code{setOn()} и читается через \code{getState()}. Мьютекс берётся внутри этих функций, поэтому
защищать данные нужно в одном месте, а не в каждой задаче.

Соберите по этой схеме отдельный скетч: кнопка и потенциометр (уроки 3 и 6), OLED (урок 7), команды
Serial (урок 2) и светодиод на ШИМ (урок 4). Экран должен обновляться независимо от нажатий.

Кнопка теперь опрашивается каждые \SI{10}{ms} независимо от того, чем занят экран: у задачи ввода
выше приоритет, и планировщик прервёт рисование, как только ей пора работать.

\begin{caution}
Переполнение стека --- частая причина перезагрузок с сообщением \code{Stack canary watchpoint triggered}.
Проверяйте запас функцией \code{uxTaskGetStackHighWaterMark(NULL)} и увеличивайте размер стека при
использовании \code{printf}, \code{String} и больших локальных массивов.
\end{caution}

\begin{trouble}
\begin{itemize}[leftmargin=*]
  \item \textbf{Перезагрузка с \code{Stack canary}} --- увеличьте размер стека задачи в \code{xTaskCreatePinnedToCore()}.
  \item \textbf{Срабатывает watchdog} --- в бесконечном цикле задачи нет \code{vTaskDelay()} или ожидания очереди.
  \item \textbf{Данные «перемешиваются»} --- общая переменная меняется без мьютекса или сразу из нескольких задач.
\end{itemize}
\end{trouble}

\begin{summary}
\begin{itemize}[leftmargin=*]
  \item Задача --- функция с бесконечным циклом; \code{vTaskDelay()} отдаёт процессор другим.
  \item Данные между задачами передаём очередью, общее состояние защищаем мьютексом.
  \item Каждое периферийное устройство (экран, Serial) лучше закрепить за одной задачей.
\end{itemize}
\end{summary}

\begin{selfcheck}
\begin{enumerate}
  \item Чем задача отличается от функции \code{loop()}?
  \item Когда нужна очередь, а когда мьютекс?
  \item Как заметить переполнение стека задачи?
\end{enumerate}
\end{selfcheck}

\begin{tasks}
\begin{enumerate}
\item \lv{1} Запустите пример с двумя задачами и убедитесь, что они работают одновременно. Измените период одной из них.
\item \lv{1} Раз в 5 секунд печатайте в Serial запас стека каждой задачи (\code{uxTaskGetStackHighWaterMark}).
\item \lv{2} Вынесите мигание встроенного RGB-светодиода в отдельную задачу. Как упростился \code{loop()}?
\item \lv{2} Передайте число из одной задачи в другую через очередь и выведите его на экран.
\item \lv{3} Переведите кнопку на прерывание (урок 3): обработчик будит задачу через \code{vTaskNotifyGiveFromISR}.
\item \lv{4} Защитите общий ресурс (например, Serial) мьютексом, а затем намеренно уберите его и найдите, когда сообщения перемешиваются.
\end{enumerate}
\end{tasks}
